\chapter{Storage e backup dei file}
\label{chap:storage}

\section{Modello di accesso SFTP-only}

Storage-Service conserva i backup degli utenti ed è l'unico componente con cui il client parli
direttamente dopo l'autenticazione. 
Una volta autenticato, infatti, l'utente diventa in grado di raggiungerlo sulla rete privata 
ed è quindi pronto per gestire i suoi backup. 

Prima di trasmetterli però, il client dell'utente, tramite Restic \cite{restic-github}, deduplica, comprime e 
cifra i backup dell'utente. Storage-Service riceve quindi una sequenza di blocchi cifrati che non è 
in grado di decifrare.
Il client, all'inizializzazione della repository Restic, prende 32 byte da \texttt{crypto/rand}, li
codifica in base64 e li usa come password del repository \cite{ramusb-code-reposecret}. 
Da quella password Restic ricava la chiave che apre il file in cui custodisce la chiave di cifratura vera, 
generata a caso alla creazione del repository \cite{restic-design}: è quest'ultima a proteggere i blocchi.

Questo approccio è stato pensato per togliere la possibilità all'utente di scegliere una password
debole. Il client salva la password del repository sul dispositivo dell'utente, con gli stessi
permessi della chiave privata SSH. La sua segretezza, come la responsabilità di copiare le chiavi in
un luogo sicuro, è lasciata all'utente. Nel caso in cui venisse persa anche una sola delle
credenziali, l'utente non potrebbe in alcun modo accedere al contenuto dei suoi backup.

Il fatto che i dati degli utenti siano cifrati, però, non significa che debbano poter essere raggiunti da 
chiunque.
L'unico modo per raggiungere i blocchi è infatti una sessione SFTP autenticata, all'interno della rete 
privata mesh. Su quella sessione, l'unica operazione possibile è trasferire file, e l'unica credenziale 
accettata è la chiave SSH \cite{openssh-sshd-config}. 

Restano fuori gli altri modi di autenticarsi e gli altri usi che SSH consente di norma: eseguire un comando, 
aprire una shell, o servirsi della connessione come tramite per raggiungere altre macchine.

Inoltre, solo gli account creati dal servizio possono presentarsi, perché il server accetta esclusivamente
nomi utente della forma \texttt{user*} e l'accesso root è esplicitamente disabilitato.

La lista predefinita di algoritmi crittografici di OpenSSH \cite{openssh-sshd-config} è volutamente ampia 
per garantire compatibilità con client che non supportano gli algoritmi moderni e più robusti, ma qui quella 
compatibilità non serve, e rischia solo di indebolire il sistema, pertanto, anche gli algoritmi 
crittografici sono elencati a mano \cite{ramusb-code-sshd-config}.

\section{Isolamento e chroot dinamico}
\label{sec:isolamento}

L'isolamento tra gli utenti si regge su un meccanismo del sistema operativo chiamato chroot, che
confina un processo dentro una sottodirectory: da quel momento quella directory è, per lui, la
radice del filesystem, e al di fuori di essa non esiste più alcun percorso che il processo possa
nominare.

Alla registrazione, Storage-Service prepara per il nuovo utente l'account POSIX che abiterà quello
spazio. L'account non ha una home directory tradizionale e la sua shell è \texttt{/usr/sbin/nologin},
un programma che si limita a rifiutare l'accesso.

Per ogni utente POSIX il cui nome corrisponde a \texttt{user*} vengono applicate anche le regole
specificate nel \autoref{lst:sshd-match}.

\Needspace{10\baselineskip}
\begin{lstlisting}[language=sshdconfig, caption={Le due direttive per utente, da
\texttt{sshd\_config} \cite{ramusb-code-sshd-config}, con commenti aggiunti.},
label={lst:sshd-match}]
Match User user*
    ChrootDirectory /storage/%u   # confina la sessione nel sottoalbero dell'utente
    ForceCommand internal-sftp    # impone SFTP al posto del comando richiesto dal client
\end{lstlisting}

Accanto all'account nasce la struttura di directory illustrata nella
\autoref{fig:filesystem-storage}. La radice del confinamento appartiene a root, mentre la
sottodirectory \texttt{data} appartiene all'utente ed è l'unico posto in cui possa scrivere. La
proprietà da parte di root è un requisito del server SSH, che
rifiuta di confinare una sessione in una directory scrivibile da chiunque altro.

Il risultato è che un utente collegato vede come radice del proprio filesystem la directory che gli
appartiene, e le directory degli altri non sono raggiungibili, né elencabili in alcun modo.

Di norma un server SSH cerca le chiavi che autorizzano una connessione in un file dentro la home
dell'utente, dove sono elencate quelle accettate per quell'account. Qui quel file non viene
consultato. 
Al suo posto, a ogni tentativo di connessione, il server esegue un binario dedicato passandogli il
nome utente che si sta presentando. Il binario gira con un'identità di sistema che non ha altri
compiti sulla macchina e chiede a Database-Vault la chiave pubblica corrente di quell'utente su un
canale mTLS.

Qualunque cosa vada storta produce lo stesso esito: il binario termina senza stampare nulla 
\cite{ramusb-code-authkeys-cmd}. Il server riceve un elenco vuoto di chiavi e nega la connessione, 
secondo il principio fail-secure \cite{owasp-fail-securely}. 
Non esiste alcun altro luogo al di fuori di Database-Vault in cui siano salvate le chiavi pubbliche SSH, 
quindi un suo malfunzionamento impedisce agli utenti di avviare le sessioni SFTP. 

\section{Backup e restore}
\label{sec:backup-restore}

Il backup avviene per intero sul dispositivo dell'utente. Il client invoca Restic indicando come
repository la sottodirectory scrivibile dentro il proprio spazio confinato, raggiunta via SFTP, e gli
impone di autenticarsi con la chiave privata generata alla registrazione
\cite{ramusb-code-restic-sftp}, che non ha mai lasciato la macchina e che nessun componente del
sistema ha mai visto.

Il percorso indicato al repository è già relativo alla radice del confinamento.

Il client risolve Storage-Service tramite MagicDNS sulla rete mesh,
e la risoluzione riesce solo finché è attivo il permesso temporaneo concesso dall'ultimo login:
scaduto quello, il nodo dell'utente non vede più il servizio e la connessione non parte nemmeno. Il
\autoref{chap:rete} descrive come quel permesso venga concesso e revocato.

La \autoref{fig:uc03} riassume la sequenza, dalla connessione SFTP fino ai dati scritti.

\begin{figure}[H]
    \centering
    \includegraphics[width=0.8\textwidth]{figures/05-usecases-sequence-uc03-backup.pdf}
    \caption{Sequenza di backup di un file (UC-03).}
    \label{fig:uc03}
\end{figure}

Il restore percorre la stessa strada in senso opposto, l'unica differenza è il comando impartito a
Restic, che scarica i blocchi richiesti e li decifra sul dispositivo dell'utente.